\section{Return lifts, physical observations, and phase coverage}
\label{sec:v83-geometry}
The clock inverse recovers absolute functions, not just their
differentials. In particular the constant in a generating function is
fixed by the zero of $T-W$; it is not chosen later by integrating a
one-form. This is essential for return time.

On a canonical domain $X\subset\R^{2h}$ with $\lambda=p\,dx$, an
exact lift is $(F,A)$ with $dA=F^*\lambda-\lambda$. Its operations are
\begin{equation}\label{eq:v83-lifts}
 (F,A)(G,B)=(F\circ G,B+A\circ G),\qquad
 (F,A)^{-1}=(F^{-1},-A\circ F^{-1}).
\end{equation}
An absolute type-I sheet $S(x,y)$ represents
\begin{equation}\label{eq:v83-generating}
 p=-S_x(x,y),\qquad p'=S_y(x,y),\qquad A(x,p)=S(x,y).
\end{equation}
It is regular when $S_{xy}$ is invertible. No minimization or convexity
of $S$ is required.

\begin{theorem}[Selective records determine a covered return lift]
\label{thm:v83-return}
Let $r,s$ be coprime positive integers. Suppose finitely many regular
endpoint gates, in fixed canonical coordinates, are given for each
of $g^r,g^s$. Their actual absolute generating sheets are the components
of Theorem~\ref{thm:v83-main} with fixed detector degree $q$. The
sheet counts, channels, actions and pairing across gates are not
supplied. Assume that these sheets are powers of one deterministic
exact lift, and that their phase graphs cover the nested domains of
a fixed word $(g^r)^u(g^s)^v$, $ur+vs=1$, including inverse-factor
sources, with fixed collars. Domains and factor order have the precise
recursive meaning in \eqref{eq:v84-word-domain}.

Then $q+3$ deadlines on each gate determine $g$ on that covered target,
including its absolute primitive. Under the margins of
Theorem~\ref{thm:v83-global}, fixed $C^{k+2}$ action bounds and uniform
inverse mixed-Jacobian and phase-map inverse bounds, the reconstruction
is locally Lipschitz from $C^{k+2}$ normalized matrix data to the $C^k$
lift. Theorem~\ref{thm:v84-compatible} gives a separate observable
replacement for the common-lift assumption on an observed invariant
domain.
\end{theorem}
\begin{proof}
The action theorem gives all unordered absolute sheets, including
coincident action values. Their derivatives give the phase graphs.
At an initial phase point, regularity solves $p=-S_x(x,y)$ locally
for $y$. Overlapping solutions have the same output and primitive
because they represent the same deterministic lift. Equality of
initial phase coordinates therefore pairs the recovered graphs,
without semantic component names. This gives $g^r,g^s$ on the required
domains, and the lift group law gives
$g=(g^r)^u(g^s)^v$. Every factor is evaluable on its nested source
by the coverage hypothesis; the primitive constants are retained.

For noisy extension fix a finite cover by smaller phase charts whose
closures lie inside the implicit-function domains, and a partition
of unity. Continue the selected reference sheet and the solution of
$p=-S_x$ using the global quotient inverse. Uniform mixed-Jacobian
bounds give $C^k$ estimates for the local phase maps and primitives
from $C^{k+2}$ sheet errors. The additional derivative bounds their
$C^{k+1}$ norms. Exact maps agree on overlaps in the fixed canonical
coordinates, so their nearby coordinate-valued extensions and
primitives may be averaged with the fixed partition. This is not
averaging unmatched sheet indices. For inverse factors continue the
local inverse branches on their own fixed source cover; the inverse
Jacobian bound gives the corresponding estimates. Compose in the
specified order on the nested collars. Product, chain and inverse
estimates give the asserted finite-word $C^k$ bound. Off the model,
the extension need not be symplectic, nor need independently extended
forward and inverse factors be exact inverses. Both properties hold
on the exact image.
\end{proof}

\begin{corollary}[An existing transverse return section]
\label{cor:v83-contact}
Let a smooth Reeb flow have a smooth transverse global section, with
first-return map $F$ and time $\tau>0$. If the actual return sheets
of two coprime counts satisfy Theorem~\ref{thm:v83-return}, selective
records recover $(F,\tau)$ and its section-preserving strict contact
suspension on the covered return domain. This determines the entire
suspension only if the entire relevant section and its return
trajectories are covered; otherwise the conclusion is the covered
suspension or flow-box germ.
\end{corollary}
\begin{proof}
Let $\lambda$ be the contact form restricted to the section. The
section flow-box pulls the contact form back to $dt+\lambda$.
Differentiating the variable-time return identity gives
$F^*\lambda-\lambda=d\tau$. Thus $(F,\tau)$ is an exact lift,
whose powers have elapsed return times as their absolute primitives.
Apply the theorem. The quotient
$(x,t)\sim(Fx,t-\tau(x))$ preserves $dt+\lambda$ on the domain
where that return identification is observed.
\end{proof}

\begin{center}\small
\begin{tabular}{@{}p{.43\textwidth}p{.48\textwidth}@{}}
\toprule
Supplied structure in Theorem~\ref{thm:v83-return}&Recovered from matrices\\
\midrule
Endpoint gates and canonical coordinates&Visible sheet count and component covering\\
Return counts and their common deterministic lift&Absolute actions and their primitives\\
Phase domains, inverse sources and collars&Canonical graphs and coordinate pairing\\
Regularity and a fixed detector family&Relative detector coefficients and weights\\
Two independent clock-invariant full-rank readouts&Both unknown channel laws\\
An existing transverse global section for a flow&The return lift and covered suspension\\
\bottomrule
\end{tabular}\end{center}
No unobserved section or continuation through an unobserved caustic is
inferred. Later we replace the common-lift premise by observable
identities and determine word domains from recovered graphs.

\subsection{An exact ambiguity outside the sampled orbit tube}
\begin{theorem}[An unsampled open phase region]
\label{thm:v83-unobserved}
Let $g=(F,A)$ be a smooth exact lift on a canonical phase domain.
For finitely many positive $n$, suppose even the full restrictions
$g^n|_{D_n}$ are known, with compact $\overline{D_n}$ and all iterates
defined. Set
\[
 K=\bigcup_n\bigcup_{j=1}^n F^j(\overline{D_n}).
\]
If $O\Subset X\setminus K$ is open and
$F^{-1}(O)\cap D_*\ne\varnothing$ for an open target $D_*$,
there are arbitrarily small smooth exact perturbations with identical
lifted observations on every $D_n$ but different values on $D_*$.
They may preserve a strict positive-roof margin on a fixed compact
neighbourhood of the target.
\end{theorem}
\begin{proof}
Choose $z_*\in F^{-1}(O)\cap D_*$. Choose a smooth Hamiltonian
with support compactly contained in $O$ and nonzero Hamiltonian vector
field at $F(z_*)$. Its sufficiently small-time flow $h_\epsilon$
stays in the canonical domain, is the identity outside $O$, and moves
$F(z_*)$ for all small nonzero $\epsilon$. Compact support allows
extension by zero to a neighbourhood of the domain, so no boundary
escape is hidden in this assertion. With
$\iota_{X_H}d\lambda=-dH$, Cartan's formula gives
\[
 B_\epsilon=\int_0^\epsilon
 (\lambda(X_H)-H)\circ h_t\dd t,
 \qquad h_\epsilon^*\lambda-\lambda=dB_\epsilon.
\]
Both the displacement and this primitive vanish outside $O$.
Put $g_\epsilon=(h_\epsilon,B_\epsilon)g$. For $z\in D_n$,
induction through $j\le n$ gives $g_\epsilon^j(z)=g^j(z)$ as lifts:
each original successive phase image lies in $K$, where the extra
lift is the identity with zero primitive. This preserves accumulated
primitives as well as phase maps. The maps differ at $z_*$. Compact
support gives convergence in every fixed finite smooth norm, so a
strict positive-roof margin on the chosen compact set persists.
\end{proof}
The result shows an unsampled-orbit freedom in the unrestricted smooth
class. It does not claim that sources illuminating additional
trajectories remain unchanged, that the construction lies in every
billiard subclass, or that every sufficient collar and phase-domain
hypothesis is separately necessary.

\subsection{Physical common clocks}
\begin{proposition}[One source and common deadlines]
\label{prop:v83-physical}
Suppose the actions on finitely many regular gates lie in $[0,H]$.
For $J=q+3$, choose $h>0$, a uniform independent release delay on
$[0,D]$, where $D=H+(J+1)h$, and deadlines $T_j=H+jh$.
Use one normalized phase source at all clocks. After the time gate,
apply unknown branchwise polynomial log-retention of degree at most
$q$ on the finite clock interval and two conditionally independent
clock-invariant readings. The model is \eqref{eq:v83-model}, with
residual and unused-delay margins at least $h$. No component audit is
needed to choose the deadlines.
\end{proposition}
\begin{proof}
The release event has probability $(T_j-W_b)/D$, with residual in
$[h,D-h]$. For a regular sheet $S_b$, the change from initial momentum
to terminal position gives successful density
\[
 \frac{\rho(x,-S_{b,x})|\det S_{b,xy}|}{D}
 \eta_b\exp\left(\sum_{\nu=1}^q\theta_{b,\nu}T_j^\nu\right)
 (T_j-W_b).
\]
All factors outside the exponential and residual are clock independent
and are absorbed in $a_b$. On the finite interval a sufficiently
small positive baseline $\eta_b$ makes retention strictly below one.
Independent conditional randomizations give the product channel laws.
A branch-blind detector or exposure only rescales the matrix.
\end{proof}
For $q=1$, a phase-dependent survival rate supplies such a detector.
Conditional independence and deadline invariance are apparatus
assumptions; copying a noisy observation does not establish them.
The detector family is prescribed, not unrestricted.

\subsection{A nonconvex return system and its persistence}
\begin{proposition}[Three sheets through an action crossing]
\label{prop:v83-shear}
There is a nonconvex exact return system with three regular sheets,
two having coincident actions along a curve, which satisfies the
log-affine hypotheses of Theorem~\ref{thm:v83-return}. These
properties persist in a sufficiently small smooth neighbourhood of
the lift on slightly smaller fixed phase and endpoint domains.
\end{proposition}
\begin{proof}
On $\R^2$ let
\[
 \phi(p)=p^3-p,\quad G(p)=\tfrac14p^4-\tfrac12p^2,\quad
 F(x,p)=(x+\phi(p),p),\quad A(p)=1+\tfrac34p^4-\tfrac12p^2.
\]
Then $F^*\lambda-\lambda=p\phi'(p)\dd p=dA$ and $A\ge11/12$.
The positive-roof quotient preserves $dt+p\,dx$ and has this return
lift. For count $n$, put $t=(y-x)/n$ and let $p_-(t),p_0(t),p_+(t)$
be the three smooth roots of $\phi(p)=t$ near $-1,0,1$. On a small
interval about zero, the absolute sheets are
\[
 S_{n,b}(x,y)=n+p_b(t)(y-x)-nG(p_b(t))=nA(p_b(t)).
\]
They satisfy $-S_{n,b,x}=S_{n,b,y}=p_b(t)$ and
$S_{n,b,xy}=-1/(n\phi'(p_b(t)))\ne0$. On the diagonal the outer
actions are both $5n/4$, the central one is $n$, and the derivative
of the outer action difference with respect to $y-x$ is $2$.
Thus the crossing is transverse and an unequal-action anchor persists
on a small compact gate. The central and outer mixed derivatives have
opposite signs.

Choose $r(p)=r_0+r_1p>0$ on the compact momentum patches, with
$r_1\ne0$, and retention $\eta e^{-r(p)(T-T_1)}$, $0<\eta<1$.
The rates separate the two outer components at their action crossing.
For one reading take
\[
 (1,e^{\beta p},e^{2\beta p})/(1+e^{\beta p}+e^{2\beta p}),
 \qquad \beta\ne0,
\]
and for the other the analogous law with a different unknown nonzero
coefficient. Independent conditional draws give two positive full-rank
channels: their columns are normalized Vandermonde columns at the
three distinct momenta. The observer is not given momenta, rates or
response coefficients. A fixed normalized smooth source positive on
the compact patches and Proposition~\ref{prop:v83-physical} give the
common clocks.

For counts two and three, momentum is unchanged and shifts are bounded
on these patches. Choose finitely many endpoint rectangles for the
successive input and output neighbourhoods of $g=g^3(g^2)^{-1}$,
then smaller target domains with collars. These supply the required
nested cover. Under a small exact smooth perturbation on the chosen
finite orbit tubes, implicit-function continuation keeps the three
regular branches of both iterates. No extra branch is included in this
restricted compact source class. Mixed-Jacobian, channel-rank, residual
and anchor margins persist, and the nonzero derivative of the outer
action difference continues the crossing curve on smaller gates.
Keep the sensors and their independent randomizations fixed as
functions of initial momentum; their rank and rate separation persist.
Collars permit the same finite gate families over smaller targets and
the strict roof bound persists. The fold points $\phi'(p)=0$ remain
outside these domains.
\end{proof}

\section{Exact boundaries of the observation theorem}
\label{sec:v83-boundaries}
\begin{proposition}[Equal actions and unrestricted detectors]
\label{prop:v83-obstructions}
If every component has one common action $w$, even arbitrarily many
projective clocks do not determine it. If the branch detector can vary
freely with deadline, small changes of unequal actions can likewise be
compensated exactly.
\end{proposition}
\begin{proof}
The common factor $T_j-w$ cancels under projectivization. In the second
case replace $W_b$ by a nearby $W'_b<T_1$ and the retention by
\[
 e'_b(T)=e_b(T)\frac{T-W_b}{T-W'_b}.
\]
All unnormalized weights remain identical. Strict residual and
retention margins preserve admissibility for small changes. The new
logarithm need not lie in the prescribed detector family, so this
does not contradict the identification theorems.
\end{proof}
\begin{proposition}[One uncalibrated reading does not suffice]
\label{prop:v83-one-view}
With one unknown reading, two unequal actions and strict positive
full-rank channels need not be identifiable at any number of clocks.
\end{proposition}
\begin{proof}
Take unit weights and zero detector coefficients. The first model has
actions $(0,1)$ and columns $(.2,.8)^{\mathsf T},(.8,.2)^{\mathsf T}$;
the second has actions $(-1,2)$ and columns
$(.4,.6)^{\mathsf T},(.6,.4)^{\mathsf T}$. For all $T>2$, both
unnormalized single-reading vectors equal
$(T-.8,T-.2)^{\mathsf T}$. The channels are positive and invertible,
but the unordered action sets differ. A common scaling gives valid
finite-clock success probabilities without changing equality.
\end{proof}
This is an obstruction to one unknown channel over the full class,
not to calibrated or otherwise supplemented observations. Enough clock
contrasts, independent readout information and an unequal-action anchor
are distinct requirements. Phase acquisition is a separate geometric
issue. The positive function-space and geometric results above identify
larger classes without removing these stated distinctions.
